% !TEX program = lualatex
% Auto-generated from YAML (v3) — 後期1: コンピュータの構成とソフトウェアの種類

\documentclass[aspectratio=43,professionalfonts,handout]{beamer}
\usepackage{beamer_template}

\newcommand{\LectureNum}{1}
\newcommand{\LectureTitle}{コンピュータの構成とソフトウェアの種類}
\newcommand{\CourseName}{情報リテラシー}
\newcommand{\TextPages}{pp.70-73}
\newcommand{\NextTopic}{処理の仕組み（CPU・命令・論理回路）}
\newcommand{\NextPages}{pp.74-77}

\begin{document}
% --- Slide 1: title ---

\begin{frame}{\CourseName\quad 第\LectureNum 回「\LectureTitle 」}
  {\small \TermName\quad \TeacherName\quad ／\quad 教科書 \TextPages}

  \vfill

  \begin{block}{今日の流れ}
    \begin{enumerate}
      \item コンピュータの基本構成
      \item ハードウェアの種類と役割
      \item 計測・制御システム
      \item ソフトウェアの種類（基本・応用）
      \item 処理の流れ
    \end{enumerate}
  \end{block}
\end{frame}

% --- Slide 2: terms ---

\begin{frame}{基本用語}
  \begin{description}[labelwidth=6em]
    \item[\kw{コンピュータ（computer）}]
      数値の記憶，演算，外部との入出力を行う機能を備えた電子機器
    \item[\kw{ハードウェア（hardware）}]
      コンピュータの筐体や部品のように，物理的に触ることができるもの
    \item[\kw{ソフトウェア（software）}]
      コンピュータを構成する物理的要素を除いた全ての無形要素の総称。一般にコンピュータプログラムを指す
    \item[\kw{基本ソフトウェア}]
      ハードウェアと応用ソフトウェアの仲介を行い，コンピュータの中核となるソフトウェア。OS（オペレーティングシステム）ともいう
    \item[\kw{応用ソフトウェア}]
      文書処理・表計算・画像処理など，特定の目的の作業に用いるソフトウェア。アプリケーションソフトウェアともいう
  \end{description}
\end{frame}

% --- Slide 3: twocol ---

\begin{frame}{ハードウェアの構成}
  \begin{columns}[T]
    \begin{column}{0.48\textwidth}
      \begin{block}{入力装置}
        \small
        \begin{itemize}
          \item キーボード・マウス・マイクロホン
          \item タッチパネル・スキャナ
          \item 外部からコンピュータへデータを取り込む
        \end{itemize}
      \end{block}
    \end{column}
    \begin{column}{0.48\textwidth}
      \begin{exampleblock}{出力装置}
        \small
        \begin{itemize}
          \item ディスプレイ・プリンタ・スピーカー
          \item コンピュータから外部へ結果を出力する
          \item 補助記憶装置（外付けHDD等）も周辺装置
        \end{itemize}
      \end{exampleblock}
    \end{column}
  \end{columns}
\end{frame}

% --- Slide 4: points ---

\begin{frame}{人間とコンピュータの対応}
  \begin{block}{ポイント}
    \begin{itemize}
      \item 入力：目・耳 → 入力装置（キーボード・マイク等）
      \item 処理：脳 → コンピュータ本体（CPU・メモリ）
      \item 出力：口・手 → 出力装置（ディスプレイ・プリンタ等）
      \item 補助記憶：手帳・ノート → 外付けHDD・USBメモリ等
    \end{itemize}
  \end{block}

  \vfill

  \begin{exampleblock}{補足}
    \begin{itemize}
      \item コンピュータは人間の行動に置き換えて考えることができる
    \end{itemize}
  \end{exampleblock}
\end{frame}

% --- Slide 5: points ---

\begin{frame}{計測・制御システム}
  \begin{block}{ポイント}
    \begin{itemize}
      \item 洗濯機・炊飯器などの家電機器にはセンサとコンピュータが組み込まれている
      \item センサ：周囲の情報を計測して電気信号に変える
      \item コンピュータ：センサからの情報を処理し，仕事を行う部分に命令する
      \item 仕事を行う部分：コンピュータからの命令に従って動作する
    \end{itemize}
  \end{block}

  \vfill

  \begin{exampleblock}{補足}
    \begin{itemize}
      \item 家電・自動車に内蔵されるこれらの制御ソフトウェアを「組み込み系ソフトウェア」という
    \end{itemize}
  \end{exampleblock}
\end{frame}

% --- Slide 6: table ---

\begin{frame}{スマートフォンの主なセンサ}
  \centering
  \small
  \begin{tabular}{lll}
    \toprule
    センサ名 & 計測対象 & 活用例 \\
    \midrule
    タッチセンサ & 接触 & 指やペンによる入力 \\
    ジャイロセンサ & 回転・傾き & 画面の縦横切り替え \\
    加速度センサ & 機器の動き・重力 & 振る動作による入力・歩数計 \\
    明るさセンサ & 明るさ & 周囲に応じた画面輝度調整 \\
    近接センサ & 近接する物体 & 通話中タッチパネルを停止 \\
    温度センサ & 温度 & 高温時に機器を停止し故障を回避 \\
    磁気センサ & 地磁気 & 方位の表示 \\
    \bottomrule
  \end{tabular}
\end{frame}

% --- Slide 7: points ---

\begin{frame}{基本ソフトウェア（OS）の主な機能}
  \begin{block}{ポイント}
    \begin{itemize}
      \item ファイル管理：ファイルの記録などを管理する
      \item メモリ管理：各種ソフトウェアが使うメモリの割り当てを管理する
      \item タスク管理：ソフトウェアの実行順序やCPUへの割り当てを管理する
      \item 入出力管理：キーボード・マウス・ディスプレイなどのハードウェアを管理する
      \item ネットワーク管理：他のコンピュータとの通信の状態を管理する
    \end{itemize}
  \end{block}

  \vfill

  \begin{exampleblock}{補足}
    \begin{itemize}
      \item OSはハードウェアの違いを吸収する働きがある。異なるPCでも同じOSなら同じように動作する
    \end{itemize}
  \end{exampleblock}
\end{frame}

% --- Slide 8: points ---

\begin{frame}{処理の流れ（ユーザ→OS→アプリ→OS→ユーザ）}
  \begin{block}{ポイント}
    \begin{itemize}
      \item ① 利用者がキーボードなどを使って情報を入力する
      \item ② 情報がハードウェアから基本ソフトウェアに渡される
      \item ③ 情報が基本ソフトウェアから応用ソフトウェアに渡される
      \item ④ 応用ソフトウェアが処理を行う
      \item ⑤ 応用ソフトウェアが出力する情報を基本ソフトウェアに渡す
      \item ⑥ 基本ソフトウェアが出力先に応じたハードウェアに情報を渡す
      \item ⑦ ディスプレイなどハードウェアからの出力が利用者に渡される
    \end{itemize}
  \end{block}
\end{frame}

% --- Slide 9: exercise ---

\begin{frame}{演習}
  {\small 各自で取り組んでください。}

  \vfill

  \begin{block}{基本問題}
    \begin{itemize}
      \item 入力装置・出力装置・補助記憶装置をそれぞれ2つずつ挙げよ
      \item 基本ソフトウェアと応用ソフトウェアの違いを説明せよ
      \item 洗濯機の計測・制御の流れをセンサ→コンピュータ→仕事を行う部分の順に説明せよ
    \end{itemize}
  \end{block}

  \vfill

  \begin{exampleblock}{応用問題（余裕がある人）}
    \begin{itemize}
      \item スマートフォンのセンサを3つ選び，それぞれARやナビゲーションへの応用例を述べよ
    \end{itemize}
  \end{exampleblock}
\end{frame}

% --- Slide 10: assignment ---

\begin{frame}{課題・次回予告}
  \begin{importantbox}[課題]
    教科書 pp.70-73 を復習すること。次週はCPUの動作原理と論理回路を学ぶ（pp.74-77）
  \end{importantbox}

  \vfill

  \begin{block}{次回}
    「\NextTopic 」（教科書 \NextPages ）
  \end{block}

  \vfill

  {\small
  質問があれば授業後またはTeamsで受け付けます。}
\end{frame}

\end{document}
